% %%%

%%%   Самарский государственный технический университет

%%%   Кафедра прикладной математики и информатики

%%%

%%%   Преамбула для статьи в журнал "Вестник Самарского государственного

%%%   технического университета. Серия: «Физико-математические науки»"

%%%   Ориентирована под MikTEx 2.4 for Win32

%%%

%%%   Хотя сам журнал собирается под GNU/Linux на дистрибутиве TeXLive



\documentclass[11pt,a4paper,twoside]{article}



\usepackage[cp1251]{inputenc}

\usepackage[T2A]{fontenc}

\usepackage[english,russian]{babel}

\usepackage{samgtubib}

\usepackage[dvips]{graphicx}

\usepackage[multiple]{footmisc}



\usepackage{samgtu}

\renewcommand{\VestnikYear}{2009} % Год издания Вестника

\renewcommand{\VestnikNumber}{2\,(19)} % Номер Вестника

\renewcommand{\VestnikMonth}{Ноябрь} % Месяц выхода Вестника

\renewcommand{\VestnikData}{01.11.09} % Дата подписания Вестника в печать

\renewcommand{\VestnikUPL}{26,64} % Усл. печ. л.

\renewcommand{\VestnikUIL}{25,73} % Уч.-изд. л.

\renewcommand{\VestnikRegNumber}{479/09} % Регистрационный номер

\renewcommand{\VestnikOrder}{994} % Номер заказа








% \usepackage{rotating_pr}


%


% \begin{document}





\renewcommand{\firstpage}{271}


\renewcommand{\lastpage}{275}





 \udk{517.956.3, 517.957}


  \msc{35L70, 35C15}





\title{Построение общего решения вырожденной системы Полмейера"--~Лунда"--~Редже}





\author{А.\,М.~Гурьева}{Г\,у\,р\,ь\,е\,в\,а~А.\,М.}





\institute{Уфимский государственный авиационный технический


университет, \\  450000, Уфа, ул. Карла Маркса, 12.}





\email{E-mails}{adel-guryeva@mail.ru}





\otherinf{Гурьева Адель Минивасимовна "--- доцент кафедры  математики; к.ф.-м.н.}





\keywords{инварианты Лапласа, уравнение лиувиллевского типа, интегралы}





\workabstract{Показано, что вырожденная система Полмейера"--~Лунда"--~Редже является


  системой  лиувиллевского типа,  получены


формулы для $x$-  и~$y$-интегралов  в~первом и~во втором порядке.


Показано, как с их помощью построить общее решение этой системы


уравнений.  }





\engtitle{Construction of General Solution of Degenerating Polmeire--Lunda--Redge System }





\engauthor{A.\,M.~Guryeva}{G\,u\,r\,y\,e\,v\,a~A.\,M.}








\enginstitute{Ufa  State Aviation  Technical  University,\\


12, Karla Marksa st., Ufa, 450000.}





\otherenginfm{Guryeva Adel Minivasimovna, Ph. D. $($Phys. \& Math.$)$,  Dept. of  Mathematics.}





\engabstract{It is demonstrated that degenerating Polmeire--Lunda--Redge system is a Liovielle-type system, formulas were obtained for $x$- and $y$-integrals at the first and the second orders. It was demonstrated how they can be used in order to construct a general solution on this equation system. }





\engkeywords{invariants of Laplace, equation of Liovielle type, integrals}





\date{14/I/2009}


\revisiondate{20/II/2009} % В окончательном виде





\maketitle





\small





 \Section[n]{Введение} В настоящее время в научной литературе


имеется огромное количество работ, посвященных гиперболическим


интегрируемым уравнениям, что и указывает на важность  данных


исследований (см., например, [1, 2]).





Определение точно интегрируемых гиперболических уравнений


\begin{equation} \label{guryeva:p1}


 u_{xy}=F(x, y, u, u_{x}, u_{y})


\end{equation}


было дано  Ж.\,Г.~Дарбу.


\begin{definition}


Уравнение (\ref{guryeva:p1}) называется {\it интегрируемым по Дарбу}, если у


него существуют нетривиальные $x$- и $y$-интегралы.


\end{definition}








Будем называть функцию $w(x, y, u, u_{x}, u_{xx}, \ldots)$, не


зависящую от $y$ на  всех решениях уравнения (\ref{guryeva:p1}),


его $x$-интегралом. Аналогично определяется $y$-интеграл.





Данная  работа  посвящена построению общего решения вырожденной


системы Полмейера"--~Лунда"--~Редже





\begin{equation} \label{guryeva:l2}


\begin{cases}


 \displaystyle u_{xy}=\frac{vu_{x}u_{y}}{uv+c}, \vspace{1mm}  \\


 \displaystyle v_{xy}=\frac{uv_{x}v_{y}}{uv+c}


\end{cases}


\end{equation}


($c$ "--- ненулевая постоянная)  с использованием $x$- и $y$-интегралов.





\Section[n]{Общее решение системы уравнений (\ref{guryeva:l2})}


Известно, что система уравнений (\ref{guryeva:l2}) является


системой лиувиллевского типа, для которой обобщённые инварианты


Лапласа первого порядка есть вырожденные матрицы, а обобщённые


инварианты Лапласа второго порядка равны нулю (см. [1,~3]).





Далее для удобства изложения материала введём следующие обозначения:


\begin{align*}


 u_{1}=u_{x}, \; u_{2}=u_{xx},  \; \ldots ,  \;  v_{1}=v_{x}, \; v_{2}=v_{xx},  \; \ldots ,  \\


\overline{u}_{1}=u_{y}, \; \overline{u}_{2}=u_{yy},  \; \ldots , \; \overline{v}_{1}=v_{y}, \; \overline {v}_{2}=v_{yy},  \;


\ldots .


\end{align*}








Построение интегралов "--- задача не из лёгких, так как порядок


интеграла заранее неизвестен. Решением этой проблемы может служить


гипотеза, высказанная А.\,В.~Жибером: порядки интегралов


и~индексы, при которых происходит  падение ранга обобщённых


инвариантов Лапласа для систем уравнений, совпадают.





Таким образом,  система уравнений (\ref{guryeva:l2}) имеет


интегралы в первом порядке:


\begin{equation}


\label{guryeva:j1}


w(x)=\frac{u_{1}v_{1}}{uv+c}, \quad


\overline{w}(y)=\frac{\overline{u}_{1}\overline{v}_{1}}{uv+c};


\end{equation}


и во втором порядке:


\begin{equation}


\label{guryeva:j2}


W(x)=\frac{u_{2}}{u_{1}}-\frac{uv_{1}}{uv+c},\quad


\overline{W}(y)=\frac{\overline{u}_{2}}{\overline{u}_{1}}-\frac{u


\overline{v}_{1}}{uv+c}.


\end{equation}











С учётом (\ref{guryeva:j1}) соотношения (\ref{guryeva:j2}) могут


быть переписаны в виде


\begin{equation} \label{guryeva:j3}


\begin{cases}{l} u_{2}-u w(x)=u_{1}W(x), \\


\overline{u}_{2}-u \overline{w}(y)=\overline{u}_{1}\overline{W}(y).


\end{cases}


\end{equation}





Относительно замены переменных $x=X(x')$, $y=Y(y')$ система уравнений (\ref{guryeva:l2})


инвариантна, поэтому в новых переменных $x'$ и $y'$ система уравнений (\ref{guryeva:j3})


принимает такой вид:


\begin{equation*}


\begin{cases}


\displaystyle\frac{\partial^{2} u}{\partial {x' }^2} \frac{1}{{X' }^2}- \frac{\partial u}{\partial {x' }}


\frac{X''}{{X' }^3 }-uw(X)= \frac{\partial u}{\partial {x' }} \frac{1}{X' }W(X), \vspace{1mm}


\\


\displaystyle\frac{\partial^{2} u}{\partial {y'} ^2} \frac{1}{{Y' }^2}-


\frac{\partial u}{\partial y' } \frac{Y''}{{Y' } ^3}-u\overline{w}(Y)=


\frac{\partial u}{\partial y' } \frac{1}{{Y' }}  \overline{W}(Y).


\end{cases}


\end{equation*}


В дальнейшем у переменных $x'$ и $y'$ штрихи опустим, тогда последняя система примет вид


\begin{equation}


\label{guryeva:j4}


\begin{cases}


\displaystyle u_{2}-u_{1} \Bigl[ \frac{ X'' }{{X' }}+W(X){X' }\Bigr]-u {X '} ^2w(X)=0, \vspace{1mm}


\\


\displaystyle \overline{u}_{2}-\overline{u}_{1}\Bigl[


\frac{Y''}{Y'}+\overline{W}(Y)Y' \Bigr]-u {Y '}^2 \overline{w}(Y)=0.


\end{cases}


\end{equation}





Выберем функции $X(x)$  и $Y(y)$ такими, что


$$


{X' }^2 w(X)=


\frac{d}{dx}\Bigl[ \frac{X'' }{X' }+W(X){X'}\Bigr], \quad


\displaystyle {Y' }^2 \overline{w}(Y)=


\frac{d}{dy}\Bigl[ \frac{Y''}{Y' }+\overline{W}(Y)Y'\Bigr].


$$


Тогда система уравнений (\ref{guryeva:j4}) переписывается так:


\begin{equation} \label{guryeva:j5}


\begin{cases}


\displaystyle u_{1}-\Bigl[ \frac{X''}{X'}+W(X){X' }\Bigr] u = \overline{A}(y),  \vspace{1mm}


\\


\displaystyle\overline{u}_{1}-\Bigl[ \frac{Y''}{Y' }+\overline{W}(Y)Y'\Bigr]u = A(x).


\end{cases}


\end{equation}





Условие совместности системы (\ref{guryeva:j5}) в виде двух


обыкновенных дифференциальных уравнений


\begin{equation*}


\overline{A}'(y)-\left[ \frac{Y^{\prime


\prime}}{Y'


}+\overline{W}(Y)Y'\right]\overline{A}(y)=c_{1}, \quad


  A'(x)-\left[ \frac{X^{\prime


\prime}}{{X' }}+W(X){X' }\right]A(x)=c_{1}


\end{equation*}


позволяет определить функции $A(x)$ и $\overline{A}(y)$:


\begin{gather*}


A(x)=c_{1}X'\Psi(X) \int \frac{dx}{X'\Psi(X)}+c_{2}X'\Psi(X),\\


\overline{A}(y)=c_{1}Y'\overline{\Psi}(Y) \int \frac{dy}{Y'\overline{\Psi}(Y)}+c_{3}Y'\overline{\Psi}(Y),


\end{gather*}


где $c_{i}= \rm const$, $i=1, 2, 3$, а $\Psi(X)$ и~$\overline{\Psi}(Y)$ "--- функции, которые удовлетворяют


таким соотношениям:


$$


\Upsilon'(X)=W(X), \quad \Psi(X)=e^{\Upsilon(X)}, \quad


\overline{\Upsilon}'(Y)=\overline{W}(Y), \quad


\overline{\Psi}(Y)=e^{\overline{\Upsilon}(Y)}.


$$





Решая последовательно обыкновенные дифференциальные уравнения в


системе (\ref{guryeva:j5}), определим $u$:


\begin{equation*}


\begin{cases}


\displaystyle u=\overline{A}(y) X'\Psi(X)


\int \frac{dx}{X'\Psi(X)}+\overline{B}(y) X'\Psi(X), \vspace{1mm}\\


\displaystyle  u=A(x) Y'\overline{\Psi}(Y)


\int \frac{d y}{Y'\overline{\Psi}(Y)}+B(y) Y'\overline{\Psi}(Y).


\end{cases}


\end{equation*}





Полученные формулы для $u$ должны совпадать. Приравнивая их, определим функции $B(x)$ и $\overline{B}(y)$:


\begin{gather*}


B(x)=c_{3} X'\Psi(X) \int \frac{dx}{X'\Psi(X)}+c_{4} X'\Psi(X), \\


\overline{B}(y)=c_{2}Y'\overline{\Psi}(Y)


\int \frac{d y}{Y'\overline{\Psi}(Y)}+c_{4}Y'\overline{\Psi}(Y),


\end{gather*}


где $c_{4} = \rm const$.


Таким образом, можно выписать формулу для нахождения функции~$u$:


\begin{equation} \label{guryeva:j6}


u=X'\Psi(X)Y'\overline{\Psi}(Y)\biggl[ c_{1}


 \!\int\!  \frac{d


x}{X'\Psi(X)} \!\int\!  \frac{d


y}{Y'\overline{\Psi}(Y)}+c_{2} \int


 \frac{d y}{Y'\overline{\Psi}(Y)}+


c_{3} \!\int\!  \frac{d


x}{X'\Psi(X)}+c_{4}  \biggr].


\end{equation}





Введём в рассмотрение функции $\Phi(x)$ и $\overline{\Phi}(y)$,  удовлетворяющие условиям


$$\Phi'(x)=\frac{1}{X'\Psi(X)},\quad  \overline{\Phi}'(y)=\frac{1}{Y'\overline{\Psi}(Y)},$$


тогда соотношение (\ref{guryeva:j6}) запишется


в компактной форме:


\begin{equation} \label{guryeva:j7}


u=\frac{1}{\Phi'(x)\overline{\Phi}'(y)}\bigl[c_{1}\Phi(x)\overline{\Phi}(y)+


c_{2}\overline{\Phi}(y)+c_{3}\Phi(x)+c_{4} \bigr].


\end{equation}





Подставляя найденное  значение $u$  в первое уравнение


(\ref{guryeva:l2}),


определяем функцию~$v$:


\begin{multline}


 \label{guryeva:j8}


v= \frac{c}{c_{2}c_{3}-c_{1}c_{4}}


\bigg[ \frac{\Phi''(x) \overline{\Phi}''(y)}{\Phi'(x)\overline{\Phi}'(y)}


\bigl( c_{1}\Phi(x)\overline{\Phi}(y)+c_{2}\overline{\Phi}(y)


+c_{3}\Phi(x)+   c_{4} \bigr)-


 \\


-  \frac{\Phi''(x) \overline{\Phi}'(y)}{\Phi'(x)}\bigl(c_{2}+c_{1}\Phi(x)\bigr)-


\frac{\overline{\Phi}''(y) \Phi' (x)}{\overline{\Phi}'(y)}\bigl(c_{3}+ c_{1}\overline{\Phi}(y)\bigr)  +  c_{1}\Phi'(x)


\overline{\Phi}'(y)\bigg].


\end{multline}





Подстановка (\ref{guryeva:j7}) и (\ref{guryeva:j8}) в последнее


уравнение системы (\ref{guryeva:l2}) даёт тождество.


Поэтому (\ref{guryeva:j7}) и (\ref{guryeva:j8}) составляют частное решение


системы уравнений (\ref{guryeva:l2}), где $\Phi(x)$ и


$\overline{\Phi}(y)$ являются произвольными функциями.








Замена


$$


u\rightarrow \alpha u, \quad v \rightarrow


 \frac{v}{\alpha}, \quad \Phi(x) \rightarrow


\varepsilon\Phi(x)+\beta, \quad \overline{\Phi}(y)\rightarrow


\lambda\overline{\Phi}(y) + \delta,


$$


где $\alpha,$ $ \beta,$ $\delta,$ $ \varepsilon,$ $ \lambda $ "--- ненулевые постоянные,


позволяет считать, что


\begin{gather*}


c_{1}=c_{4}=1 \quad \mbox{и} \quad c_{2}=c_{3}=0, \quad \mbox{


когда} \quad c_{1}\neq 0; \\


c_{2}=c_{3}=1 \quad \mbox{и}\quad c_{1}=c_{4}=0, \quad \mbox{


когда} \quad c_{1}= 0.


\end{gather*}





Тогда формулы (\ref{guryeva:j7}) и (\ref{guryeva:j8}) при $c_{1}\neq 0$  примут вид











\begin{equation} \label{guryeva:j9}


\begin{aligned}


&u= \frac{1}{\Phi'(x)\overline{\Phi}'(y)}\bigl[\Phi(x)\overline{\Phi}(y)


+1 \bigr],


\\


&v=-c\bigg[


\frac{\Phi''(x) \overline{\Phi}''(y)}{\Phi'(x)\overline{\Phi}'(y)} \bigl(\Phi(x)\overline{\Phi}(y)+1 \bigr)


- \frac{\Phi''(x) \overline{\Phi}' (y)}{\Phi'(x)}\Phi(x) +


\\


& \hspace{4.7cm}


+\Phi'(x) \overline{\Phi}'(y)


- \frac{\overline{\Phi}''(y) \Phi' (x)}{\overline{\Phi}'(y)}\overline{\Phi}(y)


\bigg],


\end{aligned}


\end{equation}


а при $c_1=0$ "---


\begin{equation} \label{guryeva:j10}


\begin{aligned}


&u= \frac{1}{\Phi'(x)\overline{\Phi}'(y)}\bigl[\overline{\Phi}(y)+\Phi(x)\bigr],


\\


&v=c\bigg[ \frac{\Phi''(x) \overline{\Phi}''(y)}{\Phi'(x)\overline{\Phi}'(y)}


\bigl( \overline{\Phi}(y)+\Phi(x)\bigr)-


\frac{\Phi''(x) \overline{\Phi}' (y)}{\Phi'(x)}-


 \frac{\overline{\Phi}''(y) \Phi' (x)}{\overline{\Phi}'(y)}\bigg].


\end{aligned}


\end{equation}








 После  преобразования $x\rightarrow F(x)$ и


$y\rightarrow\overline{F}(y)$ соотношения (\ref{guryeva:j9}) и


(\ref{guryeva:j10}) переписываются в виде


\begin{equation}


\label{guryeva:j11}


\begin{aligned}


& u= \frac{F'(x)\overline{F}'(y)}{\Phi'(x)\overline{\Phi}'(y)}\bigl[\Phi(x)\overline{\Phi}(y)+1 \bigr],


\\


&v=-c\bigg\{


\frac{F'(x)\overline{F}'(y)}{\Phi'(x)\overline{\Phi}'(y)}


\biggl[\Phi''(x) \frac{1}{{F' }^{2}(x)}-\Phi' (x) \frac{F''(x)}{{F '}^{3}(x)} \biggr]\times


\\


&\qquad \times


\biggl[\overline{\Phi}''(y) \frac{1}{{\overline{F}'}^{2}(y)}-


\overline{\Phi}'(y) \frac{\overline{F}''(y)}{{\overline{F}'}^{3}(y)}\biggr]\bigl( \Phi(x)\overline{\Phi}(y)+1 \bigr) -


\\


&\qquad -


\biggl[\Phi''(x)  \frac{1}{{F' }^{2}(x)}-


\Phi'(x) \frac{F''(x)}{{F'}^{3}(x)}\biggr] \frac{F'(x)\overline{\Phi}'(y)}{\Phi'(x)\overline{F}'(y)}\Phi(x) +


\\


&\qquad +  \frac{\Phi'(x) \overline{\Phi}'(y)}{F'(x)\overline{F}'(y)}  -


\biggl[


\overline{\Phi}''(y) \frac{1}{{\overline{F}' }^{2}(y)}-


\overline{\Phi}'(y) \frac{\overline{F}''(y)}{{\overline{F}'}^{3}(y)}\biggr]


 \frac{ \Phi^{\prime


}(x)\overline{F}'(y)}{\overline{\Phi}'(y)F'(x)}\overline{\Phi}(y)


\bigg\}


\end{aligned}


\end{equation}


при $c_{1}\neq 0$;


\begin{equation}


\label{guryeva:j12}


\begin{aligned}


& u=


 \frac{F'(x)\overline{F}'(y)}{\Phi'(x)\overline{\Phi}'(y)}


\bigl[\overline{\Phi}(y)+\Phi(x)\bigr],


\\


& v=c\bigg\{


 \frac{F'(x)\overline{F}'(y)}{\Phi'(x)\overline{\Phi}'(y)}


\biggl[


\Phi''(x)\frac{1}{{F'}^{2}(x)} -\Phi' (x)\frac{F''(x)}{{F'}^{3}(x)}\biggr]


% \times


% \\


% \times


\biggl[\overline{\Phi}''(y)\frac{1}{{\overline{F}'}^{2}(y)}-


\\


&\qquad - \overline{\Phi}' (y)\frac{\overline{F}''(y)}{{\overline{F}'}^{3}(y)}\biggr]


\bigl( \overline{\Phi}(y)+\Phi(x)\bigr)-


\biggl[ \Phi''(x) \frac{1}{{F'}^{2}(x)}-\Phi'(x) \frac{F''(x)}{{F' }^{3}(x)}\biggr] \times


\\


&\qquad \times  \frac{


 F'(x)\overline{\Phi}'(y)}{\Phi'(x)\overline{F}'(y)}-


\biggl[


\overline{\Phi}''(y) \frac{1}{{\overline{F}'}^{2}(y)}-


\overline{\Phi}'(y) \frac{\overline{F}''(y)}{{\overline{F}'}^{3}(y)}\biggr]


 \frac{\Phi'(x)\overline{F}'(y)}{\overline{\Phi}'(y)F'(x)}\bigg\}


\end{aligned}


\end{equation}


при $c_{1}=0$, соответственно.





Таким образом, общее решение системы уравнений (\ref{guryeva:l2})


может быть описано как формулой (\ref{guryeva:j11}), так и


формулой (\ref{guryeva:j12}), поскольку они сводятся  друг к другу


заменой функции $ \Phi(x)$ на $ \frac{1}{T(x)}$ и


выбором функции $ F(x) $ такой, чтобы


$F'(x)=- \frac{S'(x)}{T(x)}.$














\begin{thebibliography}{9}


 \RBibitem{[1.]}


\by Жибер~А.\,В., Соколов~В.\,В.


\paper Точно интегрируемые гиперболические уравнения лиувиллевского типа


\jour  Успехи математ. наук


\yr 2001


\vol 56


\issue 1(337)


\pages 63--106





 \RBibitem{[2.]}


\by Жибер~А.\,В., Соколов~В.\,В., Старцев~С.\,Я.


\paper О нелинейных гиперболических уравнениях,


 интегрируемых по Дарбу


\jour Докл. РАН


\yr 1995


\vol 343


\issue 6


\pages 746--748








 \RBibitem{[3.]}


\by Гурьева~А.\,М.


\paper Системы уравнений  $u_{xy}=\varphi(u, v, u_{x}, u_{y}),$


$ v_{xy}=\phi(u, v, v_{x}, v_{y}) $ с~нулевыми обобщёнными


инвариантами Лапласа второго порядка


\inbook Международн. Уфимская зимняя  шк.-конф. по математ. и физике для студент.,


аспирант. и молодых учёных


\bookinfo  Т. 1: Математика


\publaddr Уфа


\publ БашГУ


\yr 2005


\pages 195--205








\end{thebibliography}





\makeengtitlem


\lfoot[\thepage]{}


\rfoot[]{\thepage}





%


% \end{document}


